\documentclass[11pt]{article}
\usepackage[margin=2.6cm]{geometry}
\usepackage{amsmath}
\usepackage{hyperref}
\usepackage{parskip}
\setlength{\emergencystretch}{3em}

\title{Next-Track Recommendation:\\ A Plain-Language Overview}
\author{}
\date{\today}

\begin{document}
\maketitle

\begin{abstract}
\noindent A less-technical companion to \texttt{next-track.pdf} --- the same
story and the same headline numbers, with fewer acronyms. For the full
statistics, confidence intervals, and run provenance, see the main document.
\end{abstract}

\section*{What we're trying to do}

Given a listening session --- the songs you've played so far, in order ---
predict the song you'll play \emph{next}, and more broadly build a continuation
that feels right: same mood, sensible flow, not just ``more of the same album.''

We test this the honest way. We take real sessions, hide the last song, let the
model rank the whole catalog, and check whether the hidden song shows up in its
top~10. The main score, \textbf{Recall@10}, is simply \emph{``how often is the
real next song in the top~10?''}

Two things make the scores look low, on purpose:

\begin{itemize}
  \item \textbf{We only count the \emph{exact} next track.} Landing a song by the
  same artist, or the same genre, or one that simply \emph{sounds} like the
  answer, gets no credit under Recall@10 (we track those separately as softer
  scores).
  \item \textbf{Cold start is the reality.} Because we split sessions by time,
  \emph{over half} the test answers are songs that never appeared in any training
  session. No model can memorize its way to those. This caps how high any score
  can go, so every result is read \emph{relative to a simple baseline}, never as
  an absolute number.
\end{itemize}

The dataset: about 7{,}150 sessions over about 19{,}400 tracks, split into a
training half and a test half by date.

\section*{The players, in plain terms}

\begin{itemize}
  \item \textbf{The sequence model (a ``GRU'').} A small neural network that
  reads the session song-by-song and predicts what the next song should
  ``sound like,'' then finds the closest catalog songs to that guess. It captures
  the \emph{trajectory} of a session, not just the last song.
  \item \textbf{The co-occurrence baseline (a ``Markov'' model).} A simple tally
  of \emph{``which song tends to follow which''} learned from training sessions.
  It's the bar everything must beat --- and it's the model the sibling app already
  ships.
  \item \textbf{The content / ``sounds-alike'' neighbor.} For each candidate, how
  close does it sound to the session so far, using an audio-and-metadata
  fingerprint of each track. No training required.
\end{itemize}

\section*{The story, campaign by campaign}

\paragraph{1. A smart sequence model, alone, only ties the simple baseline.}
The GRU on its own matched the co-occurrence tally on Recall@10 and actually did
\emph{worse} on ranking quality. Predicting a next-song ``sound'' and hopping to
it is a harder route to the \emph{exact} next track than just remembering what
usually follows what.

\paragraph{2. Making the network fancier didn't help.}
Deeper, wider, bidirectional variants --- none beat the baseline. Architecture
was a dead end. The one lever that mattered for the network was its \emph{training
objective} (a contrastive setup was about $2.3\times$ better than the naive one),
not its size or shape.

\paragraph{3. The breakthrough: blend the two.}
Combining the sequence model with the co-occurrence baseline jumped Recall@10 by
\textbf{about $+61\%$} over the baseline. Why? The two models get \emph{different}
songs right --- one knows session flow, the other knows raw adjacency --- so
together they roughly double the hit rate. This blend became the first champion.

\paragraph{4. More blending didn't help.}
Adding a second neural model, or replacing the simple weighted blend with a
learned ``combiner,'' made things \emph{worse}, not better. Piling on more of the
same \emph{kind} of signal is redundant. The lesson: gains come from a
\textbf{genuinely new signal}, not from fancier mixing.

\paragraph{5. The one genuinely new signal was ``sounds-alike'' content.}
Adding the content neighbor as a third ingredient was the only thing that
actually improved the blend --- because it ranks \emph{different} near-misses
than the other two. That gave a three-part blend.

\paragraph{6. How we fingerprint each song matters --- and simpler won.}
Each song is turned into a compact numeric ``fingerprint.'' We compared a fancy
neural compression against a plain statistical one (PCA) at several sizes. The
plain one won at every reasonable size, and quality peaked at \textbf{192 numbers
per song} (past that it started adding noise). Intuition: the sequence model
finds the plain, well-organized fingerprint easier to aim at.

\paragraph{7. The current champion.}
The best model on record is a three-part blend where the sequence and content
legs both operate in a \textbf{learned ``sounds-like map''} --- a space tuned so
that ``close'' means ``likely to come next.'' It lands the exact next track in
the top~10 about \textbf{21\% of the time}, the right \emph{artist} about 33\%,
and the right \emph{genre} about 40\% --- far more useful than the 21\% exact
number alone suggests. Crucially, it also \textbf{reaches songs it never saw in
training} (the cold-start majority), which the other legs can't. It's been
hardened on a second, independent time-split and is \textbf{deployed and callable}
as a live model, not just a spreadsheet row.

\paragraph{8. ``Best at exact'' isn't ``best at vibe.''}
We added a continuous ``sounds-alike'' score that gives partial credit for
landing something that \emph{feels} like the answer. Interesting twist: the model
that wins on exact hits is \textbf{not} the one whose misses sound closest to the
truth. Different goals, different winners --- so we keep exact-recall as the
official score and treat ``sounds-alike'' as a diagnostic.

\paragraph{9. Why sessions can feel ``too album-eager,'' and the knob for it.}
A common complaint is that top blends parrot the current artist or album. We
measured it: the eagerness comes almost entirely from the \textbf{co-occurrence
leg}, not the neural net (the neural models are actually the \emph{least} eager).
And, against a popular hunch, the GRU and its longer-memory cousin (the LSTM) are
statistically tied --- so \emph{``the GRU wins because it's more eager''} is
\textbf{not} true. We also found a re-ranking knob that trades a \emph{tiny}
amount of exact accuracy for noticeably more variety and mood-coherence ---
essentially free at a light setting. It's available but off by default; the
official score stays exact-recall.

\section*{Bottom line}

\begin{itemize}
  \item A sequence model \textbf{alone} doesn't beat a simple
  ``what-follows-what'' tally.
  \item \textbf{Blending complementary signals} is where the wins are --- the
  two-part blend was $+61\%$, and a ``sounds-alike'' content leg added the only
  further gain.
  \item \textbf{Simpler song fingerprints} (linear, 192 numbers) beat fancier
  ones.
  \item The champion is a deployed three-part blend that's strong on exact hits,
  even stronger on getting the artist and genre right, and --- uniquely --- works
  on songs it never saw before.
  \item The ``album-eager'' feeling is real but comes from one specific leg, and
  there's a dial to soften it when we want more adventurous sessions.
\end{itemize}

\end{document}
